\section{The Modern SI and Physical Constants}\label{sec:modern-si}

High-fidelity physical simulation engines require precise numerical representations of physical constants and quantities. 
These representations must accurately reflect the true values of the constants as established by international standards.

\subsection{Defined and Experimental Constants}
Following the 2019 revision of the International System of Units (SI)~\cite{BIPM2019SIBrochure,Newell2019NIST330}, 
the definitions moved from artifact-based definitions to exact numerical definitions based on fundamental constants of nature.

Constants can be categorized into two types: 
\begin{enumerate}
\item \textbf{Defined Constants:} These are constants whose numerical values are fixed by definition in the SI standard. 
    They have zero measurement uncertainty. 
    These include $c$, $h$, $e$, $k_B$, $N_A$, and others.
    These values are hardcoded as exact literals to the maximum precision supported by the data type.
    Many other constants, such as $\mu_0$ and $\epsilon_0$, are derived from these fixed constants 
    and are also treated as exact values.
\item \textbf{Experimental Constants:} These constants are determined through empirical measurement 
    and carry experimental uncertainties\cite{Mohr2016CODATA}. 
    These include the gravitational constant ($G$), the fine-structure constant ($\alpha$), 
    and particle/object masses (e.g., electron mass $m_e$, hydrogen mass $m_H$, solar mass $M_\odot$). 
    While CODATA provides recommended values along with their uncertainties, 
    they are treated as floating-point values subject to experimental updates.
\end{enumerate}

\subsection{Simulation Engine Constants}\label{subsec:engine-constants}

To ensure reproducibility, determinism, and maximum numerical precision, 
the simulation engine initializes its internal physical constants using exact 2019 SI definitions\cite{Newell2019NIST330, BIPM2019SIBrochure} 
and standard CODATA evaluated values\cite{Mohr2016CODATA} stored as 64-bit double-precision (or higher) floating-point numbers.
The exact values can be referenced in the source code under \path{dependencies/gpp/src/core/units/units.cppm}.
